\slajdid{10-s007}
\begin{frame}[label=10-s007]
\gusttekst\setlength{\abovedisplayskip}{3pt}\setlength{\belowdisplayskip}{3pt}\setlength{\abovedisplayshortskip}{2pt}\setlength{\belowdisplayshortskip}{2pt}Kako da zaključimo u kojem režimu u kolu radi tranzistor videćemo kroz primere na logičkim kolima. Ali osnova je
\begin{enumerate}
\item Pretpostavimo režim rada tranzistora
\item Zamenimo tranzistor modelom za pretpostavljeni režim rada
\item Izračunamo parametre u kolu, napone i struje.
\item Na osnovu odnosa parametra (na primer da li je zaista $\beta_F I_B>I_C$) proverimo ispravnost naše pretpostavke.
\item Ako zaključimo da je pretpostavka pogrešna vraćamo se na tačku 1.
\end{enumerate}
Cilj „gomile“ zadataka iz elektronika jeste da u analizi lakše u startu prepoznamo u kojim režimima rade tranzistori, da ih u sintezi pravilno iskoristimo itd. Prethodno opisani postupak će voditi rešenju čak i ako više puta pogrešimo sa početnom pretpostavkom. Samo je pitanje u kojem vremenskom roku ćemo „rešiti zadatak“.
\end{frame}
